\documentclass[a4paper,openright,12pt]{book}
\usepackage[utf8]{inputenc}
\usepackage[spanish,es-noshorthands]{babel}
\usepackage{amsmath}
\usepackage{amsfonts}
\usepackage{amssymb}
\usepackage{xcolor}
\usepackage{tikz}
\usepackage{natbib}
\usetikzlibrary{arrows.meta,positioning}
\newtheorem{defi}{Definición}[section]
\newtheorem{cor}{Corolario}[section]
\newtheorem{lem}{Lema}[section]
\newtheorem{teo}{Teorema}[section]
\usepackage{ragged2e}
\usepackage{graphicx}
\usepackage[left=2.5cm,right=2.5cm,top=2.5cm,bottom=2.5cm]{geometry}
\setlength{\parindent}{0mm}

\title{Avances}
\author{jos256 }
\date{April 2021}

\begin{document}
	\textbf{\LARGE{El modelo SIR}} \hfill \break
	
	Dentro del área de la epidemiología, uno de los modelos 	
	matemáticos más usado es el modelo SIR. El cual fue propuesto 	
	por W.O Kermack y A.G Mckendrick (1927), esté modelo nos permite
	caracterizar como evoluciona una epidemia que se propaga 
	en los individuos de una cierta población (N) en un periodo de tiempo 	    (t). El modelo SIR esta formado por un sistema de tres ecuaciones 	
	diferenciales y ademas divide a la población (N) en tres estados 		
	
		\begin{itemize}
			\item{$S \left(t\right)$: Representa el número de personas 					susceptibles, personas sanas que al entrar en contacto en 	
			la enfermedad pueden resultar infectados, en función del	
			tiempo.}
			
			\item{$I \left(t\right)$: Representa el número de personas	
			infectadas, este grupo puede infectar a las personas
			susceptibles, en función del tiempo.}
			
			\item{$R \left(t\right)$: Representa el número de personas
			recuperadas, este grupo se ha recuperado de la enfermedad 
			y se han vuelto inmunes o han muerto, en función del tiempo. }
		\end{itemize}
	
	Por lo que la suma de esto tres grupos corresponde a la población   		total, esto es: 
		\begin{center}
		$N = S \left(t\right)+ I\left(t\right) + R \left(t\right)$.
		\end{center}
	El modelo SIR considera a la población constante pues no cuenta	
	los nacimientos y las muertes que suceden durante la enfermedad,
	en este modelo se considera que la forma en que se infectan las 	  		personas es por contacto directo entre las personas. Así mismo los	
	individuos del grupo $I \left(t\right)$ que logran recuperarse de 	
	la enfermedad obtienen inmunidad o perecen, esto los hace inmunes 	
	y pasan al grupo $R \left(t\right).$\nocite{allman2004mathematical}
	\hfill \break
	
	\textbf{\Large{Formulando el modelo}} \hfill \break
	
	La tasa de infección de este modelo estará dada por 	
	el producto del grupo de los susceptibles y el de los infectados, así 	
	entonces tendríamos el producto de $S \left(t\right) \cdot I \left(t		\right)$, pero al realizar esto tendríamos que todos estarían 	
	contagiados, para solucionarlo tenemos que multiplicar por un 	
	parámetro $\beta$, para así obtener $-\beta \cdot S \left(t\right) 			\cdot I \left(t	\right)$, esto ultimo pasa ya que la cantidad de 			personas de este grupo va decreciendo y se van agregando al de los 	
	infectados, a dicha $\beta$ se le conoce como tasa de contagio
	\hfill \break
	
	Para hablar de una tasa de recuperación tengamos en cuenta que los	
	individuos padecerán la enfermedad durante un tiempo determinado	
	hasta que se recupere y adquiera la inmunidad o muera.
	para esto multiplicaremos por un parámetro $\gamma$, así 	
	entonces la tasa de recuperación estará dada por $\gamma \cdot I 	    	\left(t\right)$ con $\gamma > 0$. \hfill \break
	
	por tanto se tiene que $\gamma = \frac{1}{d} $, donde 'd' es el número 		de días que pasa infectado el individuo.\hfill \break
	 
	podemos representar el sistema de ecuaciones que describe el 				sistema SIR mediante los compartimientos S,I,R y los flujos de entrada 		y salida como se muestra en el diagrama :
	
	\newpage
		
	\begin{center}
	
	\begin{tikzpicture}[node distance=2cm, auto, >=Latex, 
    every node/.append style={align=center},int/.style={draw, minimum 			size=2cm}]

   	\node [int](S) {$S$};
   	\node [int, right=of S](I) {$I$};
   	\node [int, right=of I](R) {$R$};
   	\coordinate[right=of I] (out);
   	\path[->] (S) edge node {$\beta I S$} (I)
   			  (I) edge node {$\gamma I$} (R);
	\end{tikzpicture}

	\end{center}
	
	
	Ahora formulando el sistema de EDO's que describe el comportamiento	
	del modelo SIR :
	
		\begin{equation} \label{eq1}
		\frac{dS}{dt} = -\beta S\left(t\right)I\left(t\right) \mathrm{ ; } 			\hspace{1cm}  S\left(0\right) = S_{0} > 0
		\end{equation}
					
		\begin{equation} \label{eq2}
		\frac{dI}{dt} = \beta S\left(t\right)I\left(t\right) - \gamma I				\left(t\right) \mathrm{ ; } \hspace{.4cm}  I\left(0\right) = I_{0}
		> 0
		\end{equation}			
		
		\begin{equation} \label{eq3}
		\frac{dR}{dt} = \gamma I\left(t\right) \mathrm{ ; } \hspace{2cm} R			\left(0\right) = R_{0} > 0
		\end{equation}
	
	Si tomamos a $R_{0} = 0$ estamos suponiendo que al inicio del 				proceso	de la propagación de la enfermedad, solo hay personas 				susceptibles e infecciosos, así pues $N = S_{0} + I_{0}$.\hfill \break
	
	Lo atractivo de este modelo es, que si conocemos los parámetros 	
	$\beta$, $\gamma $ y las condiciones iniciales para $S_{0}$ e $I_{0}$
	podremos darnos una aproximación de cuando la infección se propagará
	a la población y conocer el comportamiento de la epidemia.\hfill \break
	
	Pero, ¿cómo saber si habrá un brote epidémico?, para poder responder	
	esta pregunta realicemos el siguiente análisis, el cual estará dividido	
	en dos casos.\hfill \break
	
	\textbf{Caso 1:} cuando $I' > 0$ en el tiempo $t_{0}$, notemos que de 		manipular la ecuación \eqref{eq2} obtenemos $I' \left(0\right) = I_{0} 		\left( \beta S_{0}-\gamma\right) > 0$, fijémonos en $\beta S_{0}-\gamma		> 0$ por lo que se tiene $S_{0} > \frac{\gamma}{\beta} $\footnote{Obs: 		a la cantidad $\frac{\gamma}{\beta}$ también la expresan como $				\frac{1}{\rho}$}. \hfill \break
	
	De $S_{0} > \frac{\gamma}{\beta}$ obtenemos $\frac{\beta S}{\gamma}>1$
	así entonces podemos concluir que habrá un brote epidémico.
	\hfill \break
	
	\textbf{Caso 2:} cuando $I' < 0$ en el tiempo $t_{0}$ por lo que	
	$I' \left(0\right) = I_{0} \left( \beta S_{0}-\gamma\right) < 0$
	por lo que no hay epidemia ya que $S' < 0$ para cualquier instante del	
	tiempo t, por lo que $S'(t) \leq S_{0}$ para $t \geq 0$.Si consideramos 	$S_{0} < \frac{\gamma}{\beta}$ siempre se cumple que $I' \left(t\right) 	= I \left( \beta S-\gamma\right) \leq 0$ por lo que $I_{0}\geq 	
	I \left(t\right)$ para $t \geq 0$. \hfill \break
	
	De $S_{0} < \frac{\gamma}{\beta}$ obtenemos $\frac{\beta S}{\gamma}<1$
	así entonces podemos concluir que no habrá brote epidémico.
	\hfill \break
	
	Al cociente $\frac{\beta S}{\gamma}$ lo denotaremos como $R_{0}$, 
	el cual es conocido como número reproductivo básico de la infección.
	\nocite{keeling2011modeling}
	
	
	\bibliographystyle{plain}
	\bibliography{referencias}
	
\end{document}